\documentclass[11pt,a4paper]{article}
\usepackage[margin=1in]{geometry}
\usepackage{amsmath,booktabs,hyperref,float}

\title{Iteration 013: Band-Specific Spatial Filters}
\author{Autoresearch Loop}
\date{April 8, 2026}

\begin{document}
\maketitle

\begin{abstract}
Splitting the 1--9\,Hz signal into sub-bands (delta: 1--4\,Hz, theta: 4--9\,Hz)
and learning separate spatial filters per band yields $r=0.343$, significantly
worse than the unified model ($r=0.376$). Splitting destroys cross-band
coherence information that the unified filter exploits.
\end{abstract}

\section{Hypothesis}
Different frequency bands may have different optimal spatial filters due to
frequency-dependent volume conduction. Per-band CF solutions could capture
these differences.

\section{Method}
Bandpass filter input into delta (1--4\,Hz) and theta (4--9\,Hz) using FIR
filters. Learn separate closed-form spatial filters per band, then sum outputs.

\section{Results}
\begin{table}[H]
\centering
\begin{tabular}{lrrr}
\toprule
Model & Mean $r$ & Std $r$ & SNR (dB) \\
\midrule
FIR + dropout (iter011) & 0.376 & 0.076 & 0.61 \\
Band-specific CF & 0.343 & 0.063 & 0.51 \\
\bottomrule
\end{tabular}
\end{table}

\section{Analysis}
A clear negative result ($-$0.033). The 1--9\,Hz band is narrow enough that
sub-band splitting reduces effective sample size and destroys phase coherence
between delta and theta, which the unified filter uses. The lower std suggests
the model is consistently worse, not just more variable. Band splitting is
contraindicated for this narrowband regime.

\section{Next}
Iteration 014: longer FIR filters with dropout and warm restarts.

\end{document}
